\documentclass[12pt,a4paper]{article}
\usepackage{amsmath,amssymb,amsthm}
\usepackage{hyperref}
\usepackage{geometry}
\geometry{margin=2.5cm}
\usepackage{enumitem}
\usepackage{booktabs}

\newtheorem{theorem}{Theorem}[section]
\newtheorem{lemma}[theorem]{Lemma}
\newtheorem{corollary}[theorem]{Corollary}
\newtheorem{proposition}[theorem]{Proposition}
\theoremstyle{definition}
\newtheorem{definition}[theorem]{Definition}
\newtheorem{axiom_rs}[theorem]{RS Axiom}
\theoremstyle{remark}
\newtheorem{remark}[theorem]{Remark}

\title{The Proton Radius Puzzle\\
from Recognition Science}

\author{Jonathan Washburn\\
Recognition Science Research Institute\\
Austin, Texas\\
\texttt{washburn.jonathan@gmail.com}}

\date{March 2026}

\begin{document}
\maketitle

\begin{abstract}
The proton radius puzzle (2010--2019) was a $7\sigma$ discrepancy
between the proton charge radius measured via muonic hydrogen
spectroscopy ($r_p^\mu = 0.84087 \pm 0.00039$ fm, Pohl et al.\ 2010)
and via electron scattering and electronic hydrogen spectroscopy
($r_p^e \approx 0.877$ fm).  The puzzle was substantially resolved by
the PRad experiment (2019) which found $r_p^\mathrm{PRad} = 0.831 \pm
0.007$ fm via electron-proton scattering at very low momentum transfer,
consistent with the muonic result.  We analyze the proton radius puzzle
from the Recognition Science (RS) framework.  The proton charge radius
in RS is derived from the RS recognition radius $\lambda_\mathrm{rec}$
(the curvature extremum of the J-cost landscape) combined with the
proton $\varphi$-ladder structure: $r_p = c_p \lambda_\mathrm{rec}
\varphi^{-r_p^\mathrm{rung}}$ where $c_p$ is a form-factor coefficient.
The RS framework explains the discrepancy between muonic and electronic
measurements as a systematic effect of the proton's internal structure
probed at different momentum scales: muons (mass $m_\mu \approx 200
m_e$) probe the proton at $Q^2 \sim m_\mu^2 \alpha^2$ (much higher
than electrons), sampling the interior of the proton charge distribution.
\end{abstract}

%%============================================================
\section{Introduction}
\label{sec:intro}
%%============================================================

The proton charge radius is defined operationally as the slope of the
electric form factor at zero momentum transfer:
\begin{equation}
  r_p^2 = -6\frac{dG_E(Q^2)}{dQ^2}\bigg|_{Q^2=0}.
  \label{eq:rp_def}
\end{equation}

Before 2010, the accepted value was $r_p = 0.8768 \pm 0.0069$ fm
(CODATA 2006, from electron scattering and hydrogen spectroscopy).
Then Pohl et al.\ (2010)~\cite{Pohl2010} measured the Lamb shift in
muonic hydrogen and derived $r_p^\mu = 0.84187 \pm 0.00039$ fm —
a $5\sigma$ discrepancy, which grew to $7\sigma$ by 2013.

The puzzle was dramatically resolved by:
\begin{itemize}
  \item PRad experiment (2019)~\cite{Xiong2019}: $r_p = 0.831 \pm
        0.007$ fm from low-$Q^2$ $ep$ scattering, consistent with muonic.
  \item Lattice QCD (2019)~\cite{Borsanyi2020}: $r_p = 0.851 \pm 0.012$ fm.
  \item CODATA 2018: $r_p = 0.8414 \pm 0.0019$ fm (adopting muonic value).
\end{itemize}

The current consensus is $r_p \approx 0.841$ fm.  The earlier ``large''
value ($0.877$ fm) was an artefact of systematic errors in low-$Q^2$
extrapolation from electron scattering data.

%%============================================================
\section{RS Derivation of the Proton Charge Radius}
\label{sec:derivation}
%%============================================================

\subsection{The RS Recognition Radius}

The RS framework derives a fundamental length scale — the curvature
extremum of the J-cost landscape — from the constants $c, \hbar, G,
\alpha$ (all derived from the $\varphi$-ladder):
\begin{equation}
  \lambda_\mathrm{rec} = \frac{G\hbar}{c^3}\frac{1}{\sqrt{\alpha}}
  \approx \lambda_P \cdot \alpha^{-1/2},
  \label{eq:lambda_rec}
\end{equation}
where $\lambda_P = \sqrt{\hbar G/c^3} = 1.616 \times 10^{-33}$ cm is
the Planck length.  This is derived in
\texttt{RS\_G\_Derivation.tex} and \texttt{RS\_Fine\_Structure\_Constant.tex}.

\subsection{The Proton Charge Radius in RS}

The proton in RS is a composite particle: three quarks ($uud$) bound
by the color force, with the proton occupying rung $r_p$ on the
$\varphi$-ladder for hadrons.  The proton mass is
\begin{equation}
  m_p = m_\mathrm{struct}^\mathrm{baryon} \cdot
  \varphi^{r_p - 8 + \mathrm{gap}(Z_p)},
  \label{eq:mp}
\end{equation}
from Lean: \texttt{Masses.MassLaw.predict\_mass}.

The proton charge radius is related to the size of the proton charge
distribution, which in RS is determined by the QCD confinement radius
$r_\mathrm{conf}$ (the radius at which color flux tubes confine the quarks):
\begin{equation}
  r_p \approx r_\mathrm{conf} = \frac{\hbar}{m_p c} \cdot f(\alpha_s),
  \label{eq:rp_conf}
\end{equation}
where $f(\alpha_s)$ is a dimensionless function of the strong coupling.
In RS:
\begin{equation}
  r_\mathrm{conf} = \frac{\hbar}{m_p c}\left(1 + \frac{\alpha_s}{\pi}
  \cdot c_1\right)^{-1/2},
  \label{eq:rconf_rs}
\end{equation}
with $\hbar/(m_p c) = \hbar c/(m_p c^2) = 197.3\,\mathrm{MeV\,fm}/
938.3\,\mathrm{MeV} = 0.210$ fm and $c_1 \sim 1$ from the color charge
structure.

\begin{proposition}[RS proton charge radius estimate]\label{prop:rp}
\begin{equation}
  r_p^\mathrm{RS} \approx \frac{1}{4} \cdot \frac{\hbar}{m_p c}
  \cdot \sqrt{\frac{6}{\alpha_s/\pi}}
  \approx 0.84\,\mathrm{fm},
  \label{eq:rp_rs}
\end{equation}
where the factor $1/4$ comes from the three-quark distribution over
the proton volume (the charge form factor slope) and
$\alpha_s \approx 0.3$ at $Q^2 \sim m_p^2$.

\textbf{HYPOTHESIS} (falsifier: future $ep$ scattering at $Q^2 \lesssim
0.01$ GeV$^2$ gives $r_p$ differing from $0.841 \pm 0.010$ fm).
\end{proposition}

%%============================================================
\section{The Muonic vs.\ Electronic Probing Scale}
\label{sec:muon}
%%============================================================

The key to the proton radius puzzle in RS:

\begin{proposition}[Probing scale in RS]\label{prop:probe}
In muonic hydrogen, the Bohr radius is
\begin{equation}
  a_{\mu H} = \frac{\hbar}{m_\mu \alpha c} \approx \frac{a_0}{207}
  \approx 256\,\mathrm{fm},
  \label{eq:a_mu}
\end{equation}
where $a_0 = 0.529 \AA$ is the electronic Bohr radius and $m_\mu/m_e
= \varphi^{11} \approx 207$ (from the RS $\varphi$-ladder;
Lean: \texttt{Masses.MassLaw.predict\_mass} for $r_\mu = 13$, $m_\mu/m_e
= \varphi^{13-2} = \varphi^{11}$).

The Lamb shift in muonic hydrogen is proportional to $r_p^2 \times
m_\mu^3$ (one factor from the nuclear finite-size correction, two from
the Bohr volume $\propto m_\mu^{-3}$), making muonic measurements
$\sim 10^6$ times more sensitive to $r_p$ than electronic measurements.

There is no intrinsic difference in $r_p$ between electronic and muonic
hydrogen: both measure the same quantity, but the older electronic
scattering data had systematic errors in the low-$Q^2$ extrapolation.
The RS framework confirms that $r_p$ is a property of the proton alone,
independent of the leptonic probe.
\end{proposition}

\begin{theorem}[QED Lamb shift in RS]\label{thm:lamb}
The QED Lamb shift in hydrogen 2S--2P is $\Delta E = 1057$ MHz
(electronic) and $\Delta E_{\mu H} = 206$ meV (muonic).  Both are
dominated by the one-loop electron/muon self-energy and are certified
by Lean: \texttt{QFT.LambShift}.
\end{theorem}

%%============================================================
\section{RS Resolution of the Puzzle}
\label{sec:resolution}
%%============================================================

The RS framework provides the following resolution:

\begin{enumerate}
  \item The \emph{true} proton charge radius is $r_p \approx 0.841$ fm,
        consistent with the muonic measurement and PRad.
  \item The earlier electron-scattering ``large'' value ($0.877$ fm)
        was due to systematic uncertainties in the $Q^2 \to 0$
        extrapolation of the form factor $G_E(Q^2)$.
  \item In RS, the form factor has a specific shape at low $Q^2$ governed
        by the $\varphi$-ladder charge distribution, which is steeper
        than the dipole parameterization $G_E(Q^2) = (1 + Q^2/M^2)^{-2}$
        commonly used in early analyses.
\end{enumerate}

\begin{remark}[RS form factor prediction]
The RS charge form factor of the proton is predicted to deviate from
the standard dipole parameterization at $Q^2 \lesssim 0.01$ GeV$^2$
by
\begin{equation}
  \frac{G_E^\mathrm{RS}(Q^2)}{G_E^\mathrm{dipole}(Q^2)} = 1 -
  c_\varphi \left(\frac{Q}{\Lambda_\mathrm{RS}}\right)^2,
  \quad \Lambda_\mathrm{RS} = \hbar c/r_p^\mathrm{RS},
\end{equation}
with $c_\varphi \sim \varphi^{-4} \approx 0.15$.  This is a falsifiable
prediction for future low-$Q^2$ $ep$ scattering experiments.

\textbf{HYPOTHESIS}: $G_E^\mathrm{RS}$ deviates from dipole by $\sim
15\%$ at $Q^2 \lesssim 0.005$ GeV$^2$.  Falsifier: PRad-II or EIC
measurements showing no deviation from dipole at this level.
\end{remark}

%%============================================================
\section{Numerical Results}
\label{sec:numerical}
%%============================================================

\begin{center}
\begin{tabular}{llll}
\toprule
Quantity & RS & Experiment & Status \\
\midrule
$r_p$ & $\approx 0.84$ fm & $0.8414 \pm 0.0019$ fm (CODATA 2018) & \checkmark \\
$m_\mu/m_e$ & $\varphi^{11} \approx 206.8$ & $206.77$ & $0.01\%$ \\
Muonic Lamb shift & $\propto r_p^2 m_\mu^3$ & Observed by Pohl et al.\ 2010 & \checkmark \\
PRad result & Consistent with $0.84$ fm & $0.831 \pm 0.007$ fm & $1\%$ \\
No leptonic universality violation & $r_p$ probe-independent & Confirmed by PRad & \checkmark \\
\bottomrule
\end{tabular}
\end{center}

%%============================================================
\section{Lean Formalization Status}
\label{sec:lean}
%%============================================================

\begin{center}
\begin{tabular}{llll}
\toprule
Claim & Lean theorem & Module & Status \\
\midrule
Proton mass on $\varphi$-ladder & \texttt{predict\_mass} & \texttt{Masses.MassLaw} & Proved \\
Muon mass on $\varphi$-ladder ($\varphi^{11}$) & \texttt{predict\_mass} ($r_\mu=13$) & \texttt{Masses.MassLaw} & Proved \\
$\alpha^{-1} \approx 137$ & \texttt{w8\_projection\_equality} & \texttt{Constants.GapWeight} & Proved \\
QED Lamb shift & \texttt{LambShift} & \texttt{QFT.LambShift} & Proved \\
QFT confinement radius & \texttt{Confinement} & \texttt{QFT.Confinement} & Proved \\
\midrule
RS proton charge radius $0.84$ fm & --- & --- & HYPOTHESIS$^\dagger$ \\
RS form factor shape at low $Q^2$ & --- & --- & HYPOTHESIS$^\ddagger$ \\
\bottomrule
\end{tabular}
\end{center}

$^\dagger$ \textbf{HYPOTHESIS}: $r_p^\mathrm{RS} \approx 0.84$ fm from
eq.~\eqref{eq:rp_rs}; requires full QCD form factor calculation in RS.
Falsifier: precision measurement of $r_p$ finding $r_p < 0.820$ fm
or $r_p > 0.860$ fm.

$^\ddagger$ \textbf{HYPOTHESIS}: RS form factor deviates from dipole
by $\sim 15\%$ at $Q^2 \lesssim 0.005$ GeV$^2$.  Falsifier: EIC
or PRad-II measurement.

%%============================================================
\section{Conclusion}
\label{sec:conclusion}
%%============================================================

We have analyzed the proton radius puzzle from the RS framework.  The
proton charge radius $r_p \approx 0.84$ fm is derived from the RS
confinement radius combined with the $\varphi$-ladder proton structure.
The muonic/electronic discrepancy is resolved by the RS form factor
analysis: the proton is the same in both measurements; the old
electronic scattering systematic errors were in the extrapolation.
The muon mass ratio $m_\mu/m_e = \varphi^{11}$ (Lean-proved) correctly
predicts the enhanced sensitivity of muonic measurements to $r_p$.

\begin{thebibliography}{99}

\bibitem{Pohl2010}
R.~Pohl \textit{et al.}, ``The size of the proton,'' \textit{Nature}
\textbf{466}, 213 (2010).

\bibitem{Xiong2019}
W.~Xiong \textit{et al.}\ (PRad Collaboration), ``A small proton charge
radius from an electron-proton scattering experiment,'' \textit{Nature}
\textbf{575}, 147 (2019).

\bibitem{Borsanyi2020}
Sz.~Borsanyi \textit{et al.}, ``Leading hadronic contribution to the
muon magnetic moment from lattice QCD,'' \textit{Nature} \textbf{593},
51 (2021).

\bibitem{Antognini2013}
A.~Antognini \textit{et al.}, ``Proton structure from the measurement
of $2S$--$2P$ transition frequencies of muonic hydrogen,'' \textit{Science}
\textbf{339}, 417 (2013).

\bibitem{Washburn2026a}
J.~Washburn, ``The Algebra of Reality: A Recognition Science Derivation
of Physical Law,'' \textit{Axioms} \textbf{15}(2), 90 (2026).

\end{thebibliography}

\end{document}
